% Section 4: System model -- DRAFT 13 Nov 2026
% Every number is taken from the dx300-sim repo (master 260de24):
%   dx300sim/dx300_params.py, matlab/lab08a/export/boom_params_lab08a.json,
%   rl/boom_env.py, rl/boom_env_real.py, rl/10c_dynamic_check.md.
% Tags: OFFICIAL = Doosan/Develon document, ESTIMATED = drawing or literature,
% ASSUMED = modelling choice. Check before submission.

\section{System Model}
\label{sec:model}

\subsection{Boom geometry}
The study considers the boom joint of a Doosan DX300 hydraulic excavator. Fig.~\ref{fig:geometry}
% TODO: add a pin sketch (A, B, C, D) or the Excavator Design pin plot
shows the four pins used: the boom foot pin $A$ (origin), the base pin $B$ of the two boom
cylinders on the upper frame, the rod pin $C$ on the boom and the boom--arm pin $D$. With the
$x$ axis forward and $y$ up, the pin coordinates at the reference pose of the drawing
($\theta_\mathrm{ref}\approx-6.3^\circ$) are
$B=(0.571,-0.695)$\,m, $C=(2.629,0.850)$\,m and $D=(6.208,-0.682)$\,m. % ESTIMATED from a drawing
The resulting boom length $|AD|=6.245$\,m agrees with the manufacturer value. % OFFICIAL
The boom angle $\theta$ is the angle of $A\!\to\!D$ from the horizontal and is limited by end
stops to $[-47.2^\circ, 60.4^\circ]$. % ESTIMATED
The two boom cylinders (bore 140\,mm, rod 95\,mm, stroke 1450\,mm) act together; % OFFICIAL
their moment arm $r(\theta)$ about $A$ follows from the triangle $ABC$ and varies between
0.51 and 0.90\,m over the joint range.

\subsection{Joint dynamics}
The boom is modelled as a single rotational degree of freedom,
\begin{equation}
J\,\ddot\theta = \tau_\mathrm{cyl} - \tau_g(\theta) - \tau_f(\theta,\dot\theta),
\label{eq:dyn}
\end{equation}
with the gravity torque of the boom and of the front attachment lumped as point masses,
\begin{equation}
\tau_g(\theta) = g\,[\,m_b L_c\cos\theta + \textstyle\sum_i m_i(x_i\cos\theta - y_i\sin\theta)\,],
\label{eq:grav}
\end{equation}
where $(x_i, y_i)$ are the positions of the arm pin, the arm and the bucket tip in the boom frame.
Cylinder friction is a Coulomb plus viscous force mapped through the moment arm,
$\tau_f = r(\theta)\,[F_c\,\mathrm{sgn}(v) + c_v v]$ with rod speed $v=r(\theta)\dot\theta$.
The controller output $u\in[-1,1]$ sets the cylinder torque as a fraction of the
pressure-limited maximum at the 330\,bar relief pressure, % OFFICIAL
$\tau_\mathrm{cmd}=u\,\tau_\mathrm{up}(\theta)$ for $u\ge0$ (cap side) and
$u\,\tau_\mathrm{down}(\theta)$ for $u<0$ (rod side).
The pump flow $Q=494$\,L/min % OFFICIAL, 2 x 247 L/min
limits the rod speed to $v\le Q/A_\mathrm{cap}=0.267$\,m/s when extending and
$v\ge -Q/A_\mathrm{ann}=-0.496$\,m/s when retracting; on these limits the delivered torque
$\tau_\mathrm{cyl}$ is the torque that holds the speed at its limit.
The actuator energy of a move is the mechanical work
$E = \int |\tau_\mathrm{cyl}\,\dot\theta|\,dt$.
Hydraulic losses are not modelled, so $E$ is a lower bound on the pump energy.

\subsection{Two plant parameter sets}
Table~\ref{tab:params} lists the two parameter sets used.
The \emph{placeholder plant} was used for the first experiments (Section~\ref{sec:exp}):
rough values, a constant torque limit and a gravity term with one lumped mass.
The \emph{real-parameter plant} takes its values from a Simscape Multibody/Fluids model of the
DX300 boom with the front attachment. In that model the front attachment carries 71--88\,\% of
the gravity torque and about 88\,\% of the inertia, which the placeholder plant does not capture.

\begin{table}[t]
\caption{Plant parameters}
\label{tab:params}
\centering
\begin{tabular}{lcc}
\hline
Parameter & Placeholder & Real-parameter \\
\hline
Inertia about $A$, $J$ [kg\,m$^2$] & $8.0\times10^{4}$ & $2.62\times10^{5}$ \\
Boom mass $m_b$ [kg] & 5000 (lumped) & 2330 \\
Front masses (arm pin / arm / tip) [kg] & -- & 340 / 1270 / 3406 \\
CoM distance $L_c$ [m] & 3.5 & 3.12 \\
Gravity torque at $\theta=0$ [kN\,m] & 172 & 379 \\
Torque limit [kN\,m] & 500 & 514--914 (up), 277--493 (down) \\
Viscous damping & $6\times10^{5}$ N\,m\,s/rad & $c_v=2500$ N\,s/m \\
Coulomb friction $F_c$ [N] & -- & 1900 \\
Flow limit & -- & 494 L/min \\
\hline
\end{tabular}
\end{table}
% Real-parameter column: m_b, Q, relief pressure, cylinder sizes OFFICIAL; front masses
% OFFICIAL, lumping ESTIMATED; tip mass ESTIMATED (bucket 1120 kg + 1.27 m^3 x 1800 kg/m^3);
% J, Lc, friction (Li 2020, scaled) and end stops ESTIMATED.

\subsection{Numerical integration}
The controller runs at $\Delta t=50$\,ms and an episode lasts 10\,s (200 steps).
Both plants use semi-implicit Euler integration with the command held over the control step.
The placeholder plant takes one step per control step. For its low-inertia robustness case
($J\times0.38$) the damping term $h\,b/J$ reaches 0.99, and the chattering observed there was a
numerical artefact: with 10 substeps it disappears and SAC returns to its baseline accuracy.
The real-parameter plant is therefore always integrated with 10 substeps ($h=5$\,ms).
% Source: Lab 08c. Check which substep setting Table 1 (08b eval) used before submission.

\subsection{Validation against Simscape}
The real-parameter plant was compared with the Simscape model.
The static holding torque agrees within about 2\,\%.
For full-command moves across the whole joint range, the boom-up trajectory agrees
closely (arrival time $-1.2$\,\%, RMS angle difference $0.74^\circ$); Simscape adds a
$\approx$2\,Hz speed oscillation from oil compressibility, which the rigid plant omits.
For boom-down, the plant is 9.5\,\% slower (maximum angle difference $11.8^\circ$), because it
limits the retraction speed to the pump flow, whereas in Simscape gravity drives the boom
faster while the rod side cavitates. Lowering moves of the real-parameter plant are therefore
conservative.
